\chapter[Discussão]{Discussão}
\label{cap:discussao}

Este capítulo interpreta os resultados de desenvolvimento do RECLin-PT à luz da
pergunta de pesquisa, isto é, do que eles implicam para as decisões de configuração do
modelo. Discute a arquitetura adotada, as estratégias investigadas e a escolha do
\emph{encoder}, e registra as limitações e as ameaças à validade. Todas as leituras
abaixo são preliminares, porque se apoiam no DEV, com uma única semente, e aguardam
a avaliação no teste. O texto distingue o que foi observado (com remissão ao
Capítulo~\ref{cap:experimentos}), o que se pode interpretar a partir disso e o que
permanece como hipótese, isto é, explicações plausíveis que não foram testadas.

\section{A arquitetura oficial: a Pair-Aware linear}
\label{sec:discussao_estrategias}

A designação da Pair-Aware linear como arquitetura oficial do RECLin-PT não decorre
de um teste confirmatório: foi registrada depois de conhecidos os resultados do DEV,
fora do protocolo pré-registrado (Seção~\ref{sec:protocolo}). Cabe, portanto,
examinar se os resultados do DEV a sustentam e até onde.

Contra o baseline, que difere dela apenas na forma de ler a representação, o padrão
se repete nos três \emph{encoders}: a Pair-Aware linear tem macro-F1 e F1 de
\texttt{associated\_with} maiores, com diferença detectada mesmo quando se reamostram
documentos inteiros, e F1 de \texttt{negation\_of} maior, mas com diferença
detectada só com o BioBERTpt-all. Esse padrão é compatível com um benefício da
leitura explícita do par, que reduz os falsos positivos de \texttt{associated\_with}
sem sacrificar o \emph{recall} de \texttt{negation\_of}. Com uma única semente, porém,
ele não pode ser separado da variação do treino, e a repetição nos três
\emph{encoders} reduz, mas não elimina, essa dúvida, porque cada combinação foi
treinada uma única vez.

Na classe-alvo, a Pair-Aware linear não se distingue do Two-Stage em nenhum
\emph{encoder}, e a diferença para o baseline calibrado, controle que o protocolo
exige para atribuir um efeito à arquitetura e não à simples recalibração da decisão,
tem limite inferior igual a zero. Os resultados, portanto, não sustentam afirmar que
a Pair-Aware linear é superior na identificação de \texttt{negation\_of}. Sustentam
que ela alcança, sem limiar escolhido no DEV, F1 da classe-alvo equivalente ao das
demais estratégias e, em cada \emph{encoder}, o maior macro-F1 entre as estratégias
sem limiar escolhido no DEV. A variante MLP não acrescentou ganho observado sobre a cabeça linear, o
que é compatível com a escolha da cabeça mais simples.

\section{O Two-Stage e o papel do limiar}
\label{sec:discussao_twostage}

O Two-Stage deixa de ser a referência do RECLin-PT e permanece como estratégia
experimental. Pela avaliação do autor, ele não apresentou no DEV ganho suficientemente
relevante para justificar a sua maior complexidade: dois classificadores treinados
separadamente e um limiar escolhido no próprio DEV. Seus resultados ilustram um compromisso que não depende do
\emph{encoder}: o estágio 2 quase não erra o tipo dos pares que recebe, e as perdas
se concentram no estágio 1 e no limiar. Os pares que o estágio 1 não encaminha viram
\texttt{no\_relation}, o que custa algumas negações e, sobretudo, relações
\texttt{associated\_with}.

O limiar amplifica esse compromisso. Com os dois BioBERTpt, o critério do protocolo,
o F1 de \texttt{negation\_of}, escolheu o extremo da grade, $\tau=0{,}95$, à frente de
$\tau=0{,}90$ por margens de $0{,}0014$ e $0{,}0038$ no F1, com perda de cerca de
$0{,}025$ no macro-F1. Com o BERTimbau-base, escolheu $0{,}85$. A diferença entre os
\emph{encoders} no Two-Stage, nas classes que não são a alvo, decorre em grande parte
dessa escolha, e não pode ser atribuída ao \emph{encoder}. Um critério baseado no
macro-F1 teria escolhido $\tau=0{,}90$ com os dois BioBERTpt e $0{,}85$ com o
BERTimbau-base (Apêndice~\ref{ap:tau}). Como o critério foi fixado antes de qualquer
resultado, ele não foi alterado, e a consequência é registrada aqui. Do mesmo modo, a
vantagem do Two-Stage sobre o baseline na classe-alvo, de até $0{,}028$, não é
distinguível de zero nem contra o baseline nem contra o baseline calibrado, mesmo
com o viés favorável de um limiar escolhido no DEV.

\section{A escolha do \emph{encoder}}
\label{sec:discussao_encoder}

Com os três \emph{encoders} treinados nas três estratégias, o pré-treinamento clínico
não produziu vantagem detectável na classe-alvo em nenhuma estratégia: as nove
comparações de F1 de \texttt{negation\_of} têm intervalos que contêm o zero, e o
\emph{encoder} com o maior F1 muda de uma estratégia para outra. Esse quadro
contrasta com a expectativa difundida, bem como com os ganhos relatados para o
reconhecimento de entidades clínicas em português \cite{schneider_2020_biobertpt}, de
que modelos de domínio se transfiram melhor.

Duas hipóteses são compatíveis com esse resultado, sem que os experimentos atuais
permitam distingui-las ou demonstrá-las. A primeira é linguística: a negação clínica
em português é fortemente lexicalizada, e o gatilho está quase sempre colado ao achado
negado, com distância mediana de um único caractere (Seção~\ref{sec:eda_distancia}).
Um sinal tão local e regular pode depender pouco de vocabulário clínico
especializado. A segunda é estrutural: a tarefa, tal como formulada, depende da
delimitação e da direção do par, fornecidas igualmente aos três \emph{encoders} pela
mesma representação de entrada. Testar essas hipóteses exigiria, por exemplo, analisar
os erros por tipo de gatilho e por distância entre os argumentos.

Fora do Two-Stage, cujas diferenças entre \emph{encoders} acompanham o limiar, as
diferenças detectadas favorecem o BioBERTpt-clin, e não na classe-alvo: na
Pair-Aware linear e no baseline, ele tem F1 de \texttt{associated\_with} maior que o
do BioBERTpt-all, e, na Pair-Aware linear, também maior que o do BERTimbau-base.
Uma hipótese compatível, não testada, é que o pré-treinamento só em narrativas
clínicas ajude a relação menos lexicalizada, mais dependente do contexto clínico. Há,
por fim, o argumento de eficiência: com F1 de \texttt{negation\_of} equivalente, o
BERTimbau-base usa cerca de 61\% dos parâmetros, o que tende a reduzir o custo de
memória e de inferência, embora esse custo não tenha sido medido neste trabalho.

Para o RECLin-PT, a evidência disponível não basta para trocar o \emph{encoder} da
Pair-Aware linear. Pela métrica primária, os três são equivalentes no DEV; pelo
macro-F1 e pelo MCC, o BioBERTpt-clin fica à frente, mas, no macro-F1, sem diferença
detectada nos dois níveis de reamostragem, e a vantagem detectada é a de
\texttt{associated\_with}, a classe mais ruidosa da anotação; pelo custo, o
BERTimbau-base é o menor. Por isso, o BioBERTpt-all, fixado no protocolo, continua
sendo o \emph{encoder} de referência, e os demais seguem como variantes
experimentais, que só substituem a referência se assim indicar a seleção da
configuração final, feita no DEV antes do congelamento.

\section{A classe \texttt{associated\_with}}

A classe \texttt{associated\_with} é a mais difícil para todas as estratégias, e a
análise dos falsos positivos (Seção~\ref{sec:resultados_assoc}) indica de onde vem
parte da dificuldade. Cerca de um quarto dos falsos positivos são pares cuja direção
inversa está anotada como associação: o modelo reconhece o vínculo entre as duas
entidades, mas não a direção escolhida pela anotação. Como quase nenhuma associação do
DEV é anotada nas duas direções, e como a associação, ao contrário da negação, não
tem uma direção linguística evidente, uma hipótese plausível é que a direção anotada
seja, em parte, arbitrária para o modelo. A precisão também cai com a distância entre
os argumentos, o que é compatível com a hipótese de que distinguir associação
genuína de mera proximidade seja ambíguo em texto clínico telegráfico, em que
entidades vizinhas frequentemente não mantêm relação anotada, como no exemplo do
Quadro~\ref{quad:exemplo_candidatos}.

Dois fatores medidos limitam, além disso, o que se pode esperar da classe. O F1 é
calculado sob o teto de \emph{recall} imposto pela janela, que deixa de fora cerca de
8\% das relações \texttt{associated\_with} do DEV, contra cerca de 3\% das de
\texttt{negation\_of} (Seção~\ref{sec:teto_recall}). E, nas seis notas anotadas duas
vezes no corpus, a concordância entre as duas anotações de \texttt{associated\_with} é
de cerca de $0{,}57$ (Seção~\ref{sec:particao}), um indício de que a própria
referência é instável nessa classe. Como o foco deste trabalho é a negação, nenhuma
decisão foi otimizada para \texttt{associated\_with}.

\section{Limitações}
\label{sec:limitacoes}

As limitações deste trabalho são agrupadas em cinco categorias: experimentais,
arquiteturais, dos dados, de escopo e de reprodutibilidade. Todas foram observadas
ou registradas ao longo do desenvolvimento, e as remissões indicam onde cada uma é
fundamentada.

\subsection{Limitações experimentais}

A limitação mais séria é a ausência de avaliação no teste. Todos os números são de
DEV, o mesmo conjunto usado para escolher épocas e limiares. Os do Two-Stage e do
baseline calibrado são otimistas por construção, e mesmo as escolhas de época
introduzem algum otimismo nas demais estratégias. A avaliação definitiva será feita no
teste, uma única vez, depois de congelada a configuração final. Pela emenda proposta,
essa avaliação se restringe ao modelo final. Assim, o teste estimará o desempenho do
RECLin-PT, mas não a sua vantagem sobre o baseline ou sobre as demais estratégias,
que continuará apoiada só no DEV. Os experimentos de aprimoramento previstos até o
congelamento, avaliados no mesmo DEV, tendem a aumentar esse otimismo.

Há uma única semente, de modo que não é possível separar o efeito de uma estratégia
ou de um \emph{encoder} da variação do treino: os intervalos de confiança medem só a
incerteza sobre os exemplos. O processo de seleção também limita a leitura dos
resultados. A designação da Pair-Aware linear como arquitetura oficial foi registrada
depois de conhecidos os resultados do DEV e não é confirmatória, e o protocolo
registrado ainda prevê as comparações principais do Two-Stage no teste, enquanto a
emenda proposta (E3) não for aprovada (Seção~\ref{sec:protocolo}). O protocolo fixa
o BioBERTpt-all, de modo que as execuções com o BioBERTpt-clin e o BERTimbau-base são
variantes de \emph{encoder} fora do pré-registro, e a Pair-Aware MLP é exploratória.
Em todas as execuções de três rótulos, a época escolhida foi a última permitida, o
que sugere que mais treino poderia mudar os valores absolutos
(Seção~\ref{sec:curvas}); o limiar $\tau$ com os dois BioBERTpt e o
\texttt{min\_freq} do léxico ficaram no extremo das respectivas grades; e os
controles do protocolo (baseline calibrado e filtro léxico) só existem com o
BioBERTpt-all.

\subsection{Limitações arquiteturais}

Os achados se limitam à formulação de três classes adotada, com apenas
\texttt{negation\_of} e \texttt{associated\_with} como relações, e à escala
\emph{base} dos modelos (Seção~\ref{sec:encoders}). Entre as estratégias
investigadas, o Two-Stage não tem otimização conjunta dos dois estágios e depende de
um limiar escolhido no DEV (Seção~\ref{sec:twostage}), e a cabeça da Pair-Aware foi
avaliada apenas nas formas linear e com uma camada oculta. O \emph{fine-tuning}
restrito, previsto no protocolo, não foi executado. Outras possibilidades
arquiteturais, como modelos de maior porte ou tipos semânticos na entrada, ficaram
deliberadamente fora deste trabalho e permanecem para investigação futura.

\subsection{Limitações dos dados}

A anotação apresenta inconsistências que afetam o treino e a avaliação. $38$
relações \texttt{negation\_of} do treino e $5$ do DEV têm a direção invertida, com a
expressão de negação como segundo argumento. Cerca de $37\%$ das entidades
\texttt{Negation} não participam de nenhuma \texttt{negation\_of} anotada, de modo
que parte dos pares \texttt{no\_relation} pode corresponder a relações omitidas, e
não a ausências reais (Seção~\ref{sec:controles_tipo}). Nas seis notas anotadas duas
vezes no corpus, a concordância entre as duas anotações (F1) é de cerca de $0{,}57$
nas entidades e nas relações \texttt{associated\_with}, contra $1{,}0$ nas $7$
relações \texttt{negation\_of} (Seção~\ref{sec:particao}). Variações na delimitação
de entidades e omissões de anotação no SemClinBr também foram relatadas por
\citeonline{ramos_2025_clinpthyb}.

Quanto ao alinhamento e à delimitação, a reconstrução dependeu de uma regra de
alinhamento dos \emph{offsets} aplicada a $421$ documentos; a auditoria não encontrou
divergência entre o trecho recortado e o texto anotado, mas as $26$ anomalias de
anotação encontradas, como um \emph{offset} inválido, foram mantidas no corpus
(Seção~\ref{sec:parsing}). Há ainda particularidades da organização dos dados: três
notas do teste são cópias, diferentes só em espaços, de notas do treino; há uma
única partição, na qual a semente $42$ deixou o DEV com poucas relações
\texttt{associated\_with} e o teste com muitas
(Tabela~\ref{tab:sensibilidade_semente}); e a janela de candidatos impõe um teto de
\emph{recall} mais baixo para \texttt{associated\_with} do que para
\texttt{negation\_of} (Seção~\ref{sec:teto_recall}). Essa janela
(\texttt{max\_gap}$=25$) foi definida a partir das distâncias das relações
\emph{gold} do corpus completo, inclusive das anotações do teste.

\subsection{Limitações de escopo}

O RECLin-PT trabalha com entidades previamente anotadas no corpus. O sistema não
realiza o reconhecimento automático de entidades integrado à extração de relações, e
o erro que uma etapa automática de reconhecimento introduziria não é considerado nas
métricas. Os resultados descrevem, portanto, a extração de relações em condições
ideais de entrada. A integração entre NER e RE permanece como possibilidade de
evolução do sistema (Capítulo~\ref{cap:conclusao}).

\subsection{Limitações de reprodutibilidade}

Persistem pontos de procedência que precisam ser resolvidos antes da avaliação no
teste (Apêndice~\ref{ap:inventario}). As execuções foram produzidas com quatro
versões distintas do código, registradas em sete revisões, enquanto o protocolo prevê
uma identidade única na conferência final. As diferenças entre as versões,
conferidas nos manifestos de código registrados por cada execução, correspondem às
mudanças documentadas da Pair-Aware, à criação da variante MLP e à inclusão dos novos
\emph{encoders} na configuração. O código da estratégia baseline é o mesmo em todas
elas. Sete execuções foram retomadas a partir de um \emph{checkpoint} e registram
apenas a sessão final de treino; os estágios do Two-Stage não registram a revisão do
código nem o campo de avaliação do teste; na composição do Two-Stage com o
BERTimbau-base, os nomes dos estágios registrados no limiar diferem dos nomes das
pastas, embora os \emph{hashes} da configuração e do modelo confiram; os
\emph{checkpoints} e os registros de console ficam fora do repositório e não foram
verificados; e todas as sessões registradas trazem um aviso de não determinismo de
uma rotina de atenção da GPU, emitido apesar de os algoritmos determinísticos terem
sido solicitados. Nenhum desses pontos altera as métricas recalculadas a partir das
predições, mas todos limitam a reprodutibilidade estrita e reforçam a necessidade de
preservar versões, configurações e registros experimentais até a avaliação final.

\section{Ameaças à validade}

Quanto à validade interna, as partições congeladas e verificadas por
\emph{hash}, a reconstrução dos dados com offsets alinhados ao texto e o
particionamento em nível de documento mitigam o vazamento de dados e os erros de
alinhamento. Duas ameaças permanecem registradas: as três notas do teste com cópia no
treino, para as quais o protocolo prevê relatar o teste também sem esses documentos
(Seção~\ref{sec:particao}), e a definição da janela de candidatos
(\texttt{max\_gap}$=25$), feita com as distâncias das relações \emph{gold} do corpus
completo, inclusive as do teste (Seção~\ref{sec:teto_recall}). Essa segunda ameaça é
limitada: a decisão envolveu apenas anotações, não predições ou métricas de modelos,
é idêntica para todos os sistemas comparados e afeta a cobertura dos candidatos, e
não a seleção de épocas, limiares ou configurações, feita só no DEV. Ainda assim, a
cobertura do teste não pode ser considerada estimada de forma independente dessa
escolha; no teste, por exemplo, todas as relações \texttt{negation\_of} anotadas
ficaram dentro da janela (Tabela~\ref{tab:teto_recall}).

Uma segunda ameaça diz respeito à independência das instâncias. O \emph{bootstrap}
previsto no protocolo trata os pares candidatos como independentes, embora pares de
um mesmo documento compartilhem texto. Para dimensionar esse efeito, as comparações
do DEV foram repetidas com o \emph{bootstrap} por documento, que produziu intervalos
em geral mais largos e, em algumas comparações, deixou de detectar diferenças
detectadas por candidato (Seção~\ref{sec:resultados_comparacoes}). As conclusões
deste capítulo se apoiam nas diferenças detectadas nos dois níveis. Quanto à
validade externa, o SemClinBr representa um conjunto específico de instituições e de
estilos de redação, e as 26 anomalias de anotação mantidas no corpus impõem um limite
ao desempenho alcançável, igual para todos os sistemas.
